%% kapa: exam=1 round=26 year=2015 period=1 form=A subject=economics subject_ko=경제원론 seq=8 number=48
%% source: Q-Net 감정평가사 기출문제 (cst00309) artlSeq=5138001 file=extracted/1008/5138001/제26회_감정평가사 1차_기출/제26회_감정평가사 1차_기출/1교시_감평_A형.pdf
%% conversion: kordoc 4.24.51 -> tools/gen1.py (auto); review=proofread-sonnet; flags=none
\begin{question}{48}
효용극대화를 추구하는 소비자 甲의 효용함수는 $U(x,y)=x+y$이다. 甲의 무차별곡선에 관한 설명으로 옳지 않은 것은? (단, 甲은 $X$재와 $Y$재만 소비하고 $x$는 $X$재의 소비량, $y$는 $Y$재의 소비량이며, $x$, $y$는 양수이다.)
\begin{choices}
  \item 원점에서 멀리 있는 무차별곡선은 원점에서 가까이 있는 무차별곡선보다 선호된다.
  \item 무차별곡선은 우하향한다.
  \item 무차별곡선들은 서로 교차하지 않는다.
  \item 동일한 무차별곡선 상에서 한계대체율은 체감한다.
  \item 무차별곡선의 기울기는 모든 소비조합(consumption bundle)에서 동일하다.
\end{choices}
\end{question}
\answer{4}
